\documentclass[10pt,a4paper]{amsart}
\usepackage[margin=19mm]{geometry}
\usepackage{amsmath,amssymb,mathtools}
\usepackage[scr=rsfs,cal=euler]{mathalfa}
\usepackage{xfrac}
\usepackage[unicode,hidelinks,bookmarks=false]{hyperref}
\newcommand{\ie}{i.e.\ }
\newcommand{\cf}{cf.\ }
\newcommand{\resp}{respectively\ }
\newcommand{\Subsection}{\subsection}
\providecommand{\nobreakdash}{\nobreak\hbox{-}}
\setlength{\emergencystretch}{2em}
\multlinegap0pt
\usepackage{amssymb}
\usepackage{mathtools}
\newtheorem{proposition}{Proposition}
\newcommand{\C}{{\mathbb{C}}}
\newcommand{\Ocal}{{\mathcal{O}}}
\title{Algebraic Mapping Class Group Rigidity}
\author{Seong Youn Kim}
\date{Source: arXiv:2508.09421v2 (CC BY 4.0)}
\begin{document}
\maketitle

\noindent\emph{Source and licence.} Algebraic Mapping Class Group Rigidity. Version arXiv:2508.09421v2 (CC BY 4.0), distributed by arXiv under the Creative Commons Attribution 4.0 International licence, http://creativecommons.org/licenses/by/4.0/, as recorded in the arXiv licence metadata for this exact version and shown on the abstract page https://arxiv.org/abs/2508.09421v2. Full licence notice: \texttt{licenses/arxiv-2508.09421-CC-BY-4.0.txt}.\par\smallskip

\noindent\textit{Editorial scope.} Counted unit: the article's proposition DimIntersection (source0.tex lines 737--743). The article's complete original proof of it is delivered with it (source0.tex lines 745--774, one \texttt{proof} environment of 30 lines). No mathematics is written, repaired or re-derived: the statement and the proof are the article's own lines, in the article's own notation, and no retained line is substituted or re-flowed.  No further source-local context is needed: every cross-reference of every retained line resolves inside the two delivered spans.  One declared adaptation: exactly 3 ancillary strings inside the retained spans are omitted (image-inclusion commands and, where the source repeats a label, the repeated label definitions) and no other line differs from the source; the figure environments, with their placement, captions and labels, are kept, and the package records the omitted strings in fidelity receipts bound to the affected bodies.  No retained line contains a citation, so the wrapper carries no bibliography and no bibliography entry.  The retained lines contain the article's own diagram code, which the wrapper typesets with the standard tikz packages; no image file is delivered and the package declares no asset dependency.  Editorial formatting only, in addition: an AMS \texttt{amsart} preamble that restates the article's own declarations and macros used by the retained lines, namely the environments \texttt{proposition} on separate counters, together with the standard packages those lines need.  The retained environments print in this document as Proposition 1 in order of appearance, whereas the article prints the same items with the numbering of its own full document.  The complete native source of the article as distributed in the arXiv e-print for this exact version is bundled unchanged as \texttt{sources/183168/source0.tex}.
\begin{proposition}\label{DimIntersection}
Let $\gamma, \delta$ be non-parallel essential simple closed curves. Then
\[
\dim \mathcal{O}_{\gamma} \cap \Ocal_{\delta} \leq \dim A_\Sigma - 2
\]
where the equality holds if and only if $i(\gamma,\delta)=0$ or $i(\gamma, \delta) = 1$ and $\Sigma = \Sigma_{1,1}$.
\end{proposition}

\begin{proof}
We may change the base to $\C$. Observe $\Ocal_\gamma \cap \Ocal_\delta$ corresponds to the image of $\C[X(\Sigma |_{(\gamma \cup \delta)})] \rightarrow \C[X(\Sigma)]$. Thus
\[
\dim \Ocal_\gamma \cap \Ocal_\delta \leq \dim \C[X(\Sigma |_{(\gamma \cup \delta)})]
\]
and it suffices to compute $\dim \C[X(\Sigma |_{(\gamma \cup \delta)})]$.

If $i(\gamma, \delta)=0$, the equality holds by the previous proposition.

If $i(\gamma, \delta)>0$, take the representatives of $\gamma, \delta$ with minimal intersection.
We will consider $\delta$ as arcs in the cut surface $\Sigma |_\gamma$. $\delta$ is divided into $i(\delta, \gamma)$ pieces of arcs connecting the boundary components. By cutting along each arc segment, the pair $(g, b)$ changes where $g$ denotes the genus and $b$ denotes the number of boundary components of the cut surface. 

By cutting along each arc, there are three possibilities. If the arc is separating, $(g, b) \rightarrow (g, b+1)$. This process will be denoted by \textit{pants cut}, a half cut of a (local) pants. If the arc is nonseparating, $(g, b) \rightarrow (g-1, b+1)$ or $(g, b) \rightarrow (g, b-1)$. These will be denoted by \textit{glued pants cut} and \textit{boundary merging}.

\begin{figure}[h]
    \centering
    
    
    
    \caption{(i) Pants cut. (ii) Glued pants cut. (iii) Boundary merging. }
\end{figure}

In any case the Euler characteristic increases by 1. If the cut pieces do not contain any annuli or polygons, then the dimension decreases by $3$. 

Note that polygons are yielded only by boundary merging, and for this one needs a precedent annulus. Thus if one has a cut polygon then the dimension drops at least $3$.

A pants cut yields an annulus if and only if one of the leg is indeed a leg of pants, which is not local. Similarly, a glued pants cut yields an annulus if and only if it is a cut of a once-punctured torus. Finally, a boundary merging yields an annulus if and only if the component is a pants.

One concludes that if $\dim \Ocal_\gamma \cap \Ocal_\delta = \dim \C[X(\Sigma)] - 2$ then $i(\gamma, \delta) = 1$ and an annulus is yielded by a boundary merging. Thus $\Sigma = \Sigma_{1, 1}$.
\end{proof}

\end{document}
